% ============================================================
%   Exercice Simplexe - Sujet 3 : Voyageur de commerce (TSP)
%   Fichier LaTeX autonome, compilable directement par :
%       pdflatex simplexe.tex
% ============================================================
\documentclass[12pt,a4paper]{article}

% --- Encodage et langue ---
\usepackage[utf8]{inputenc}
\usepackage[T1]{fontenc}
\usepackage{lmodern}
\usepackage[french]{babel}

% --- Mathematiques ---
\usepackage{amsmath}
\usepackage{amssymb}
\usepackage{amsfonts}
\usepackage{mathtools}

% --- Mise en page ---
\usepackage{geometry}
\geometry{margin=2.5cm}
\usepackage{parskip}

% ============================================================
%   COMMANDE \Eq : signe egal renforce (visible dans tout lecteur
%   PDF, y compris VS Code PDF Viewer). Police texte + Large + bold.
% ============================================================
\newcommand{\Eq}{\mathrel{\text{\Large\bfseries =}}}

% ============================================================
%   COMMANDES \Leq et \Geq : inegalites <= et >= en version
%   inclinee (plus visible et bien distincte de < et >).
% ============================================================
\newcommand{\Leq}{\leqslant}
\newcommand{\Geq}{\geqslant}

% --- Equation encadree avec \Eq ---
\newcommand{\BoxedEq}[1]{%
  \begin{equation*}
    \boxed{\,\displaystyle #1\,}
  \end{equation*}%
}

\begin{document}

\section*{Exercice Simplexe : Probleme du voyageur de commerce (TSP)}

\subsection*{Contexte}
Le probleme du voyageur de commerce consiste a trouver un circuit optimal permettant de visiter plusieurs points. Pour appliquer la methode du simplexe, on considere ici une version lineaire simplifiee avec deux types de trajets.

\subsection*{Variables de decision}
\[
\begin{aligned}
X &\Eq \text{nombre de trajets effectues sur la route 1}, \\
Y &\Eq \text{nombre de trajets effectues sur la route 2}.
\end{aligned}
\]

\subsection*{Fonction objectif}
L'objectif est de maximiser le gain obtenu par le choix des trajets :

\BoxedEq{Z_{\mathrm{max}} \Eq 25\,X + 20\,Y}

\subsection*{Contraintes}
\[
\left\{
\begin{aligned}
X + Y    &\Leq 30  \\
4X + 2Y  &\Leq 100 \\
2X + 3Y  &\Leq 90  \\
X, Y     &\Geq 0
\end{aligned}
\right.
\]

\subsection*{Forme standard}
On ajoute les variables d'ecart \(S_1\), \(S_2\) et \(S_3\) :
\[
\begin{aligned}
X + Y + S_1            &\Eq 30  \\
4X + 2Y + S_2          &\Eq 100 \\
2X + 3Y + S_3          &\Eq 90  \\
Z - 25X - 20Y          &\Eq 0
\end{aligned}
\]

\subsection*{Resolution}
Les sommets realisables principaux sont :
\[
(0,0), \quad (25,0), \quad (0,30), \quad (20,10)
\]

Calcul de la fonction objectif a chaque sommet :
\[
\begin{aligned}
Z(0,0)    &\Eq 0 \\
Z(25,0)   &\Eq 25(25) + 20(0)  \Eq 625 \\
Z(0,30)   &\Eq 25(0)  + 20(30) \Eq 600 \\
Z(20,10)  &\Eq 25(20) + 20(10) \Eq 700
\end{aligned}
\]

\subsection*{Solution optimale}

\BoxedEq{X \Eq 20 \quad\text{et}\quad Y \Eq 10}

\BoxedEq{Z_{\mathrm{max}} \Eq 700}

\subsection*{Interpretation}
La meilleure combinaison consiste a effectuer 20 trajets sur la route 1 et 10 trajets sur la route 2. Le gain maximal obtenu est de 700 unites.

\end{document}
